\frontchapter{I. What This Book Sets Out to Do}
\section*{1. Intended Readers}
\begin{itemize}
\item Young people with middle school reading skills who are curious about the human world.
\item Readers who think history is just dates, literature just a central message, philosophy just quotations, religion just superstition or doctrine, and art just a question of whether something looks good.
\item Adults who want a broad map of the humanities but do not want to begin with difficult terminology and weighty classics.
\end{itemize}
\section*{2. Central Aims}
This book does more than bind language, history, politics, and art lessons together. It follows questions running throughout human history:

\begin{itemize}
\item Why do humans need language, and how does language change thought and society?
\item Why do people tell stories, and how can literature invent events yet still tell us truths?
\item How do people confront death, suffering, chance, and what they cannot understand?
\item What knowledge deserves our trust? How do we distinguish argument, rhetoric, evidence, and authority?
\item Why do societies develop laws, states, hierarchies, markets, and revolutions?
\item Why do humans make art? What makes something beautiful, or makes it art?
\item How did history happen, and on what grounds do we claim to know the past?
\end{itemize}
The book follows one chronological thread:

\textbf{Human evolution $\rightarrow$ language and ritual $\rightarrow$ farming and cities $\rightarrow$ writing and states $\rightarrow$ classical civilizations $\rightarrow$ world religions and philosophical traditions $\rightarrow$ exchanges across continents $\rightarrow$ maritime expansion $\rightarrow$ science, Enlightenment, and revolution $\rightarrow$ industrialization, nation-states, and imperialism $\rightarrow$ world wars and decolonization $\rightarrow$ globalization, the digital age, and artificial intelligence.}

It also uses six lenses: \textbf{language, literature, religion, philosophy, history, and aesthetics}.

\section*{3. After Reading, You Should Be Able To}
\begin{itemize}
\item Read a basic global historical timeline without equating world history with European history or a succession of dynasties.
\item Analyze a view's premises, evidence, reasoning, and hidden assumptions.
\item Distinguish factual statements, causal explanations, value judgments, and declarations of position.
\item Recognize that language has rules yet constantly changes, and that ``correct language'' also involves context and power.
\item Attend to narrators, viewpoints, structure, tone, imagery, and what remains unsaid when reading stories.
\item Understand religions respectfully and comparatively, without rushing to praise, dismiss, or conflate them.
\item Understand why freedom, justice, happiness, truth, and beauty have no one-sentence answers, yet cannot be answered arbitrarily.
\item Ask about sources, context, and alternative explanations when encountering historical narratives, news, short videos, and online arguments.
\end{itemize}
\section*{4. What This Book Is Not}
\begin{itemize}
\item A guide to memorizing famous quotations, dates, intellectual schools, and book titles.
\item A textbook promoting or opposing a particular religion, political system, or philosophical position.
\item An account reducing complex traditions to convenient slogans such as ``the East values intuition, the West reason.''
\item A history of inevitable progress in one direction, from ``backwardness'' to ``modernity.''
\item A reduction of literature to standard answers or aesthetics to rankings of taste.
\item An exercise in false balance that pretends every position has equal evidence and equal plausibility.
\end{itemize}
\section*{5. Suggested Scope}
The material lends itself to three volumes and twelve parts:

\begin{itemize}
\item \textbf{Volume I, People Who Speak and Tell Stories}: Humanities methods, language, literature, prehistory, and early civilizations.
\item \textbf{Volume II, The Sacred, Reason, and Order}: The classical world, religion, philosophy, and global connections in the medieval period.
\item \textbf{Volume III, The Modern World and the Question of Beauty}: Early modern transformations, modern history, aesthetics, the digital age, and shared human problems.
\end{itemize}
For a single-volume publication, close readings of classic texts, reconstructions of philosophical arguments, historical debates, and art theory could be placed in a companion volume for deeper reading.
